\chapter{El Puente Cuántico: Clifford+T y la Compatibilidad con QPU}

\section{Introducción a la Computación Cuántica y Unitariedad}
El marco de POLYDIM se basa fuertemente en transformaciones ortogonales puras, que en un espacio vectorial complejo se generalizan naturalmente a transformaciones unitarias ($U^\dagger U = I$). Esto abre la puerta a una simulación directa sobre unidades de procesamiento cuántico (QPU). Los estados latentes se asimilan a registros de qubits entrelazados, y la evolución del enjambre es descrita por puertas cuánticas.

\section{El Grupo de Clifford y la Puerta T}
El grupo de Clifford $\mathcal{C}_n$ está generado por las puertas de Hadamard ($H$), Fase ($S$) y CNOT. Por el \emph{Teorema de Gottesman-Knill}, cualquier circuito compuesto exclusivamente por compuertas de Clifford y mediciones en la base computacional puede ser simulado en una computadora clásica en tiempo polinomial $O(n^2)$. En POLYDIM, las transformaciones puramente simétricas pueden formularse como circuitos de Clifford. 

Para alcanzar la universalidad cuántica completa, se requiere una compuerta no-Clifford, típicamente la compuerta $T = \text{diag}(1, e^{i\pi/4})$. El grupo generado por Clifford+T es denso en $SU(2^n)$, permitiendo aproximar cualquier operador unitario.

\section{Síntesis Exacta y el Algoritmo de Ross-Selinger}
El algoritmo de Ross-Selinger proporciona una síntesis óptima de rotaciones $Z$ utilizando compuertas Clifford+T.
\begin{theorem}[Ross-Selinger]
Cualquier rotación de un solo qubit $R_z(\theta)$ puede ser aproximada hasta una precisión $\epsilon$ con un recuento de puertas $T$ acotado por $T\text{-count} \le 3 \log_2(1/\epsilon) + O(\log \log 1/\epsilon)$.
\end{theorem}
La síntesis exacta ocurre en el anillo $\mathbb{Z}[1/\sqrt{2}, i]$, que es el ambiente algebraico natural para las isometrías ortogonales en el espacio euclidiano que subyace a la deformación $SU_q(2)$ de POLYDIM.

\section{Clifford Twirling y Compilación Aleatorizada}
A medida que los agentes en el MAS aplican rotaciones sucesivas, el error de coma flotante clásico se acumula de forma coherente en $O(N \epsilon)$. En el hardware cuántico análogo, el ruido coherente destruye la fidelidad del estado.
El \emph{Clifford Twirling} transforma el ruido coherente arbitrario en un canal de Pauli incoherente:
\begin{equation}
    \mathcal{E}(\rho) = \sum_{k} p_k P_k \rho P_k^\dagger
\end{equation}
Esta compilación aleatorizada mitiga la acumulación coherente, forzando al error a crecer como $\sqrt{N}$ (paseo aleatorio en S^{D-1}), un paralelismo directo con las estrategias de estabilización de POLYDIM (Caterpillar to Butterfly).

\section{Corrección de Errores GKP y Flotantes Subnormales}
El código GKP (Gottesman-Kitaev-Preskill) codifica un qubit lógico en los modos continuos de un oscilador armónico, ofreciendo protección contra pequeños desplazamientos en el espacio de fases (hasta $\sqrt{\pi}/2$). En la arquitectura clásica de POLYDIM, el análogo directo de estos desplazamientos infinitesimales inmanejables son los números de coma flotante subnormales. La aplicación repetida del operador límite asintótico actúa análogamente a la recuperación por corrección de síndrome de GKP, limpiando las oscilaciones bajo la barrera de precisión.

\section{El Invariante de Bargmann-Pancharatnam}
Cuando un agente experimenta una serie de transformaciones latentes formando un bucle cerrado en el espacio de fases de $S^{D-1}$, adquiere una fase geométrica global irremovible, expresada por el invariante de Bargmann-Pancharatnam:
\begin{equation}
    \Phi_{BP} = \arg\left( \langle \psi_0 | \psi_1 \rangle \langle \psi_1 | \psi_2 \rangle \dots \langle \psi_n | \psi_0 \rangle \right)
\end{equation}
En POLYDIM, monitorear $\Phi_{BP}$ en simulaciones locales certifica si la secuencia de rotores de Clifford ha completado exitosamente un bucle geodésico o si el proceso ha divergido hacia trayectorias disipativas, sin requerir trazabilidad explícita del tensor a lo largo del tiempo.

\section{Hardware Físico: Mapeo de Rotores}
La correspondencia entre los rotores en el álgebra de Clifford $C\ell(D, 0)$ y la sintaxis QPU permite portar experimentos de POLYDIM directamente a hardware como Cerebras WSE-3, IBM Quantum, Rigetti y Google Sycamore.

\section{El Teorema de Solovay-Kitaev y la Eficiencia de la Síntesis}

El Teorema de Solovay-Kitaev establece que el número de compuertas universales
necesario para aproximar cualquier unitaria $U \in SU(2)$ con precisión $\varepsilon$ es:
\begin{equation}
n_{\text{gates}} = \mathcal{O}\!\left(\log^c\!\!\left(\frac{1}{\varepsilon}\right)\right), \quad c \approx 3.97
\end{equation}

Para la tolerancia certificada del kernel V762 $\varepsilon = 4.44 \times 10^{-16}$:
\begin{equation}
n_{\text{gates}} \approx \left(\log_2(2.25 \times 10^{15})\right)^{3.97} \approx 52^{3.97} \approx 5.8 \times 10^6 \text{ compuertas}
\end{equation}

Esto es viable en hardware cuántico con miles de qubits lógicos (proyección 2028--2030).

\section{Magic States y la Destilación T}

Las compuertas $T$ son el recurso escaso en computación cuántica tolerante a fallos.
El protocolo de destilación de Bravyi-Kitaev produce 1 magic state de alta fidelidad a partir de:
\begin{equation}
n_{\text{noisy}} = 15^k \text{ copias ruidosas, con error } \varepsilon_{\text{out}} = \varepsilon_{\text{in}}^{3^k}
\end{equation}

Para $\varepsilon_{\text{in}} = 10^{-3}$ y alcanzar $\varepsilon_{\text{target}} = 10^{-15}$: nivel $k = 12$.

\section{Escalabilidad Cuántica: Qubits Requeridos para $D = 10^6$}

Para representar un vector en $S^{D-1}$ con $D = 10^6$:
\begin{equation}
n_{\text{qubits\_lógicos}} = \lceil \log_2 D \rceil = \lceil \log_2 10^6 \rceil = 20 \text{ qubits}
\end{equation}

Con código superficial de distancia $d = 15$ ($\approx 1000$ qubits físicos por lógico):
\begin{equation}
n_{\text{qubits\_físicos}} = 20 \times 1000 = 20{,}000 \text{ qubits físicos}
\end{equation}

IBM Quantum (2025) supera $1000+$ qubits físicos. El hardware para POLYDIM cuántico
está proyectado para 2028--2030 con sistemas de $>10{,}000$ qubits de alta fidelidad.

\section{Síntesis Cuántica Discreta Clifford+T en Rust (Arquitectura V772)}
\label{sec:discrete_clifford_t}

En arquitecturas cuánticas tolerantes a fallos (FTQC), las rotaciones analógicas continuas $R_y(\theta)$ son físicamente irrealizables sin introducir ruido de calibración continuo. La arquitectura V772 integra un compilador cuántico discreto nativo en Rust con interfaz C-ABI.

\begin{theorem}[Descomposición Canónica de Rotaciones Monocúbit]
\label{thm:clifford_canonical}
Toda rotación arbitraria de Pauli $R_y(\theta)$ y $R_x(\theta)$ en la esfera de Bloch se sintetiza de forma exacta a partir de rotaciones axiales $R_z(\theta)$ conjugadas por elementos del grupo de Clifford:
\begin{align}
R_y(\theta) &= H \cdot R_z(\theta) \cdot H \\
R_x(\theta) &= H \cdot S \cdot R_z(\theta) \cdot S^\dagger \cdot H
\end{align}
donde $H$ es la compuerta de Hadamard y $S = T^2$ es la compuerta de fase.
\end{theorem}

\begin{proof}
Recordando las identidades de Pauli: $H X H = Z$ y $H Z H = X$. Para el eje $Y$, dado que $S^\dagger X S = Y$, se tiene $H S^\dagger X S H = H Y H$. Conjugando las exponenciales matriciales correspondientes:
\begin{align}
H e^{-i \frac{\theta}{2} Z} H &= e^{-i \frac{\theta}{2} H Z H} = e^{-i \frac{\theta}{2} X} = R_x(\theta) \\
H e^{-i \frac{\theta}{2} X} H &= e^{-i \frac{\theta}{2} H X H} = e^{-i \frac{\theta}{2} Z} = R_z(\theta)
\end{align}
Aplicando la transformación de similitud con Hadamard a la rotación $R_z$, se obtiene directamente $R_y(\theta) = H R_z(\theta) H$. \qed
\end{proof}

\subsection{Compilación Diádica y Optimización del $T$-Count}
La rotación axial $R_z(\theta)$ se aproxima mediante descomposición en fracciones diádicas utilizando el algoritmo de \emph{GridSynth} (Ross \& Selinger) sobre la base universal $\{H, S, T, CX\}$. El compilador genera una secuencia discreta de compuertas:
\begin{equation}
R_z(\theta) \approx U_m U_{m-1} \cdots U_1, \quad U_k \in \{H, S, T\}
\end{equation}
con un costo de compuertas $T$ acotado por $n_T = 3 \log_2(1/\varepsilon) + \mathcal{O}(\log(\log(1/\varepsilon)))$, logrando una síntesis determinista con error acotado en silicio local (Suite 5 PASS).
